\clearpage
\SourceLeaf{87}%

\begin{center}
\large LAGRANGES LEHRSATZ,\\ AUF MÖGLICHST KLARE WEISE HERGELEITET.
\end{center}

\Needspace{8\baselineskip}
\begin{center}\textls[90]{I. Spezieller Fall.}\end{center}

\noindent Grundgleichung: \quad $x = t + u x^{\lambda},$

\noindent Gesucht: \quad $x^n = t^n + T_1 u + T_2 uu + T_3 u^3 + \text{u.\,s.\,w.}$

\noindent Hier sind $T_1$, $T_2$, $T_3$ usw. Funktionen von $t$, die in $\Theta_1$, $\Theta_2$, $\Theta_3$ usw. übergehen mögen, wenn man anstelle von $n$ überall dort, wo es in ihnen vorkommt, $n+\lambda$ schreibt. Es gilt also offenbar:

\[ x^{n+\lambda} = t^{n+\lambda} + \Theta_1 u + \Theta_2 uu + \Theta_3 u^3 + \text{u.\,s.\,w.} \]

Allgemein gilt für eine beliebige Funktion $Fx$ von $x$

\[ \frac{dFx}{du} = x^{\lambda} \mathord{.} \frac{dFx}{dt}, \]

\noindent also

\[ \frac{dx^n}{du} = \frac{x^{\lambda} dx^n}{dt} = \frac{n}{n+\lambda} \frac{dx^{n+\lambda}}{dt} \]

\noindent und folglich:

\[\begin{aligned}
T_1 + 2 T_2 u + 3 T_3 uu + \text{u.\,s.\,w.}
&= \frac{n}{n+\lambda}\Bigl\{(n+\lambda)t^{n+\lambda-1}+\frac{d\Theta_1}{dt}u\\
&\quad+\frac{d\Theta_2}{dt}uu+\frac{d\Theta_3}{dt}u^3+\text{u.\,s.\,w.}\Bigr\}.
\end{aligned}\]

\noindent Es gilt also $T_1 = n t^{n+\lambda-1}$, und daraus ergibt sich, wenn man $n+\lambda$ anstelle von $n$ schreibt,

\[ \Theta_1 = (n+\lambda)\, t^{n+2\lambda-1}; \]
